\documentclass[11pt]{article}
\usepackage[margin=1in]{geometry}
\usepackage{amsmath,amssymb,bm,booktabs,array}
\usepackage[T1]{fontenc}
\usepackage{lmodern}
\usepackage[hidelinks]{hyperref}
\setlength{\parindent}{0pt}
\setlength{\parskip}{0.65em}
\newcommand{\R}{\mathbb{R}}
\newcommand{\skew}[1]{[#1]_{\times}}
\begin{document}

\begin{center}
{\LARGE\bfseries MIT VNAV 2023 -- Lab 2 Individual Solutions}\\[0.6em]
{\large Wenqian Zhang \quad Student ID: 24020036087}
\end{center}
\hrule

\section*{Deliverable 1: Nodes, Topics, and Launch Logic}
The application nodes are \texttt{/av1broadcaster},
\texttt{/av2broadcaster}, \texttt{/plots\_publisher\_node}, and
\texttt{/rviz}. The two static broadcasters publish on \texttt{/tf\_static}.
The plots node obtains transforms through its tf2 listener (which receives
\texttt{/tf} and \texttt{/tf\_static}) and publishes the drone meshes and
trajectory trails on \texttt{/visuals}. RViz subscribes to these transform and
visualization topics.

If the \texttt{static} argument is omitted, it takes its default value
\texttt{false}. Thus the group guarded by \texttt{if="\$(arg static)"} is
skipped, while the node guarded by \texttt{unless="\$(arg static)"} is
launched. The latter is \texttt{frames\_publisher\_node}, so the vehicle
frames move and are published continuously on \texttt{/tf}.

\section*{Deliverables 2 and 3: ROS/tf Implementation}
The submitted \texttt{two\_drones\_pkg} contains the completed source. The
frames publisher uses one timestamp for both transforms, the elapsed ROS time
\(t\), translations
\[
{}^{w}\bm{o}_1=(\cos t,\sin t,0)^T,\qquad
{}^{w}\bm{o}_2=(\sin t,0,\cos 2t)^T,
\]
and the yaw quaternion \(q=\operatorname{RPY}(0,0,t)\) for AV1. The plots
publisher obtains the newest available transform using
\texttt{lookupTransform(ref\_frame,dest\_frame,ros::Time(0))}.

\section*{Deliverable 4: Relative Trajectory}
\subsection*{1. AV2 in the world frame}
Because \(x=\sin t\), \(y=0\), and \(z=\cos 2t=1-2\sin^2 t\), AV2 follows
the parabola
\[
y=0,\qquad z=1-2x^2,\qquad -1\le x\le 1.
\]

\subsection*{2. AV2 expressed in the AV1 frame}
AV1 has rotation \({}^{w}R_1=R_z(t)\). Therefore
\[
{}^{w}T_1=\begin{bmatrix}R_z(t)&{}^{w}\bm{o}_1\\0&1\end{bmatrix},\qquad
{}^{w}T_2=\begin{bmatrix}I&{}^{w}\bm{o}_2\\0&1\end{bmatrix},
\]
and
\[
{}^{1}T_2=({}^{w}T_1)^{-1}{}^{w}T_2
=\begin{bmatrix}
R_z(-t)&R_z(-t)({}^{w}\bm{o}_2-{}^{w}\bm{o}_1)\\0&1
\end{bmatrix}.
\]
The relative position is consequently
\[
{}^{1}\bm{o}_2=
\begin{bmatrix}
\sin t\cos t-1\\-\sin^2t\\\cos2t
\end{bmatrix}
=\begin{bmatrix}
\tfrac12\sin2t-1\\\tfrac12(\cos2t-1)\\\cos2t
\end{bmatrix}.
\]

\subsection*{3. Plane containing the relative trajectory}
From the last two components,
\(z=\cos2t=2y+1\). Hence the trajectory lies in
\[
\Pi:\ z-2y-1=0,
\]
whose normal can be chosen as \(\bm{n}=(0,-2,1)^T\).

\subsection*{4. A convenient coordinate frame on the plane}
The ellipse center is
\(\bm{p}_c=(-1,-\tfrac12,0)^T\). Choose the orthonormal basis
\[
\bm{e}_{x'}=\begin{bmatrix}1\\0\\0\end{bmatrix},\quad
\bm{e}_{y'}=\frac1{\sqrt5}\begin{bmatrix}0\\1\\2\end{bmatrix},\quad
\bm{e}_{z'}=\frac1{\sqrt5}\begin{bmatrix}0\\-2\\1\end{bmatrix}.
\]
After subtracting \(\bm{p}_c\) and projecting onto this basis,
\[
x'=\tfrac12\sin2t,\qquad
y'=\tfrac{\sqrt5}{2}\cos2t,\qquad z'=0.
\]

\subsection*{5. Shape and semiaxes}
Eliminating \(t\) gives
\[
\frac{x'^2}{(1/2)^2}+\frac{y'^2}{(\sqrt5/2)^2}=1.
\]
Thus the relative trajectory is an ellipse centered at \(\bm{p}_c\), with
semiaxis lengths \(1/2\) and \(\sqrt5/2\).

\section*{Deliverable 5: Quaternion Identities}
Use scalar-last unit quaternions \(q=[\bm{q}^T,q_4]^T\) and define
\[
\Omega_1(q)=\begin{bmatrix}
q_4I+\skew{\bm q}&\bm q\\-\bm q^T&q_4
\end{bmatrix},\qquad
\Omega_2(q)=\begin{bmatrix}
q_4I-\skew{\bm q}&\bm q\\-\bm q^T&q_4
\end{bmatrix}.
\]
These matrices represent left and right quaternion multiplication. Since
\(\skew{\bm q}^T=-\skew{\bm q}\) and
\(\skew{\bm q}^2=\bm q\bm q^T-\|\bm q\|^2I\), direct block
multiplication yields
\[
\Omega_i(q)^T\Omega_i(q)
=(\|\bm q\|^2+q_4^2)I_4=I_4,\qquad i\in\{1,2\}.
\]
Geometrically, multiplication by a unit quaternion preserves quaternion norm,
so its matrix representation must be orthogonal.

Moreover,
\[
\Omega_1(q)^Tq=\Omega_2(q)^Tq=
\begin{bmatrix}\bm0\\1\end{bmatrix}=\bm e_4.
\]
This is the matrix form of \(q^{-1}\otimes q=q\otimes q^{-1}=1\).

Finally, left multiplication and right multiplication commute as operators:
for every quaternion \(z\), associativity gives
\[
x\otimes(z\otimes y)=(x\otimes z)\otimes y.
\]
Therefore
\[
\Omega_1(x)\Omega_2(y)=\Omega_2(y)\Omega_1(x).
\]
Replacing \(y\) by its conjugate (whose multiplication matrix is the
transpose for a unit quaternion) also gives
\[
\Omega_1(x)\Omega_2(y)^T=\Omega_2(y)^T\Omega_1(x).
\]

\section*{Optional Deliverable 6: Extrinsic and Intrinsic Rotations}
Rotations about fixed/world axes premultiply the current attitude. Applying
\(R_0\), then \(R_1\), then \(R_2\) therefore produces
\[
R_{\mathrm{ext}}=R_2R_1R_0.
\]
Rotations about moving/body axes postmultiply the current attitude. Applying
the same rotations in reverse order, \(R_2\), then \(R_1\), then \(R_0\),
produces
\[
R_{\mathrm{int}}=I\,R_2R_1R_0=R_{\mathrm{ext}}.
\]
Hence an extrinsic sequence is equivalent to the reverse-order intrinsic
sequence.

\end{document}

